\section{Custom Floats}\label{ch6sec:custom-floats}

The template already provides custom float-like environments for side-by-side figures and shaded text blocks. If you need additional float types---for example, ``algorithm'' or ``protocol'' floats---you can define them with the \texttt{newfloat} package.

\begin{figure}[h]
\centering
\subfloat[][]{\begin{tikzpicture}
\draw[thick, primary, fill=primary!10] (0,0) rectangle (3,1.5);
\node[align=center, font=\small] at (1.5,0.75) {Algorithm};
\end{tikzpicture}}
\hspace{1em}
\subfloat[][]{\begin{tikzpicture}
\draw[thick, primary, fill=primary!10] (0,0) rectangle (3,1.5);
\node[align=center, font=\small] at (1.5,0.75) {Protocol};
\end{tikzpicture}}
\mycaption{Custom floats}{Placeholder figures illustrating custom float types.}
\label{ch6fig:setup}
\end{figure}

\subsection{Defining a New Float}

A new float type is declared with \verb|\DeclareFloatingEnvironment|. For example, to create an \texttt{algorithm} float, write:
\begin{verbatim}
\DeclareFloatingEnvironment[name=Algorithm]{algorithm}
\end{verbatim}
You can then create an \texttt{algorithm} environment, add a \verb|\caption|, and include it in a list of algorithms. Remember to add the new list to the table of contents if required.

\begin{table}[h]
\caption{Custom float types in the template}\label{ch6tab:demo-nav}
\centering
\begin{tabular}{l l}
\toprule
\thead{Float type} & \thead{Defined via} \\
\midrule
Figure & Standard \texttt{figure} environment \\
Table & Standard \texttt{table} environment \\
Listing & \texttt{longlisting} environment \\
Side-by-side figures & \texttt{sidebysidefigures} environment \\
\bottomrule
\end{tabular}
\end{table}

\subsection{Integration with the Template}

Custom floats should follow the same styling conventions as figures and tables. Use the \texttt{mdframed} package to add a coloured frame if desired, and use \verb|\mycaption| for captions. Table~\ref{ch6tab:demo-nav} summarises the float types available in the template.

\begin{figure}[ht]
\centering
\begin{tikzpicture}
\draw[thick, primary, fill=primary!10] (0,0) rectangle (6,2.5);
\node[align=center] at (3,1.25) {\textbf{Custom Float Demo}\\ Define, caption, list.};
\end{tikzpicture}
\mycaption{Qualitative comparison}{A placeholder qualitative comparison figure.}\label{ch6fig:qualitative}
\end{figure}
